\documentclass[12pt,a4paper]{article}

\usepackage{amsmath,amssymb,amsfonts}
\usepackage{geometry}
\usepackage{hyperref}
\usepackage{booktabs}
\usepackage{physics}
\usepackage{xcolor}
\usepackage{graphicx}
\usepackage{natbib}

\geometry{margin=2.5cm}

\hypersetup{
  colorlinks=true,
  linkcolor=blue,
  citecolor=blue,
  urlcolor=blue
}

% ─── Macros SSEE ────────────────────────────────────────────────────────────
\newcommand{\phiG}{\varphi}
\newcommand{\KAL}{\mathrm{KAL}_0}
\newcommand{\KALeff}{\mathrm{KAL}_{\rm eff}}
\newcommand{\MIRA}{\mathcal{M}}
\newcommand{\AURA}{\mathcal{A}}
\newcommand{\Mpl}{M_{\mathrm{Pl}}}
\newcommand{\Omde}{\Omega_{\mathrm{DE}}}
\newcommand{\Omm}{\Omega_{m}}
\newcommand{\half}{\tfrac{1}{2}}
\newcommand{\Hzero}{H_0}
\newcommand{\bc}{\beta_c}
\newcommand{\fsc}{f_{\mathrm{screen}}}
\newcommand{\alphaK}{\alpha_K}
\renewcommand{\Omega}{\varOmega}
\newcommand{\Mcuv}{M^4_{\rm UV}}
\newcommand{\rhoc}{\rho_{\rm crit}}
% ─────────────────────────────────────────────────────────────────────────────

\title{\textbf{UV Completion of SSEE Dark Energy:\\
       $M^4 = 45\alpha^2\rho_{\rm crit} = 5\phiG^8\rho_{\rm crit}$\\
       and a Conditional Near-Exact Hubble Tension Resolution}}

\author{Mike Edison Almeida Vallejo\\
        \small\texttt{mike.almeida1721@gmail.com}}

\date{May 2026 --- Paper 10 of the SSEE series}

\begin{document}
\maketitle

\begin{abstract}
Within the Structural Self-Energy Expansion (SSEE-V3.6), the infrared (IR) kinetic
Lagrangian $K(X) = X/\KAL$ predicts a local Hubble constant
$H_0^{\rm IR} = 72.86$ km/s/Mpc ($0.17\sigma$ from SH0ES).
We extend the EFT to include the leading higher-derivative operator
$K(X) = X/\KAL + X^2/M^4$ and identify the unique algebraic value of the
UV cutoff $M$ such that the theory remains connected to the inflationary
$\alpha$-attractor of Paper~1.
The $\alpha$-attractor parameter $\alpha = \phiG^4/3$, which fixes the
inflationary tensor-to-scalar ratio $r = 12\alpha/N^2 \approx 0.00813$
(a pre-committed LiteBIRD target), also determines
\begin{equation*}
  \boxed{M^4 \;=\; 45\,\alpha^2\,\rhoc \;=\; 5\,\phiG^8\,\rhoc
         \;\approx\; 234.9\,\rhoc}\,,
\end{equation*}
where the identity $45\alpha^2 = 5\phiG^8$ follows exactly from the definition
$\alpha = \phiG^4/3$ (not a numerical coincidence, but an algebraic consequence).
With this cutoff, the full Bellini--Sawicki parameter becomes
$\alphaK^{\rm full} = 0.41691$ ($+3.37\%$ over the IR value), giving
$H_0^{\rm local} = 73.040$ km/s/Mpc, consistent with the SH0ES
central value ($73.04 \pm 1.04$ km/s/Mpc; Riess et~al.\ 2022) to within the
rounding precision of the published SH0ES result.
\textbf{This agreement is conditional on Postulate~C.1 (UV-IR $\varphi/\pi$
separability; §\ref{sec:theorem})} and is not independent of SH0ES: the
postulate was formulated with SH0ES as the target.  It should therefore be
regarded as a self-consistency check, not an independent prediction.
The same $\alpha$ that governs inflation at $z \sim 10^{60}$ thus fixes the
dark-energy UV scale today, providing a falsifiable cross-epoch consistency test:
if LiteBIRD measures $r \neq 0.00813$, both $\alpha$ and $M$ change, and the
predicted $H_0^{\rm local}$ changes accordingly.  The IR prediction
$H_0^{\rm IR} = 72.86$ km/s/Mpc ($0.17\sigma$ SH0ES) is independent of
this postulate and stands on its own.
\end{abstract}

\tableofcontents
\newpage

% ═══════════════════════════════════════════════════════════════════════════════
\section{Introduction}
% ═══════════════════════════════════════════════════════════════════════════════

The Hubble tension---the $\sim4$--$5\sigma$ discrepancy between CMB-inferred and
local distance-ladder measurements of $H_0$---has resisted resolution through
conventional extensions of $\Lambda$CDM.
Paper~9 of this series~\citep{SSEE9} showed that within SSEE-V3.6, the EFT
kineticity parameter $\alphaK = 0.4033$ (Papers 7--8) and the lensing factor
$\MIRA = \AURA/2$ combine to produce an algebraic screening fraction
\begin{equation}
  \fsc \;=\; \frac{\alphaK}{3\,\MIRA}
       \;=\; \frac{\pi - \phiG}{(\phiG+\pi)^2}
       \;=\; 0.06725\,,
\end{equation}
which corrects the algebraic Hubble constant $H_0^{\rm alg} = 3(\phiG+\pi)^2
= 67.96$ km/s/Mpc to
\begin{equation}
  H_0^{\rm IR} \;=\; \frac{H_0^{\rm alg}}{1 - \fsc} \;=\; 72.86 \text{ km/s/Mpc}
  \quad (0.17\sigma \text{ from SH0ES}).
\end{equation}

That IR result uses the kinetic Lagrangian $K(X) = X/\KAL$ (the leading term),
and the UV cutoff $M \to \infty$ is implicit.
This paper asks: \emph{what is $M$ physically?}

The leading higher-derivative correction is $K(X) = X/\KAL + X^2/M^4$.
We show that the unique algebraically consistent value of $M$ is
\begin{equation}
  M^4 \;=\; 45\,\alpha^2\,\rhoc \;=\; 5\,\phiG^8\,\rhoc\,,
  \label{eq:M4main}
\end{equation}
where $\alpha = \phiG^4/3$ is the inflationary $\alpha$-attractor parameter of
Paper~1, and the equality $45\alpha^2 = 5\phiG^8$ is an exact algebraic identity
following from the definition of $\alpha$.
With this value, the corrected Hubble constant is $H_0^{\rm UV} = 73.040$
km/s/Mpc, consistent with the SH0ES central value $73.04 \pm 1.04$ km/s/Mpc
to within the rounding precision of the published result (conditional on
Postulate~C.1; see §\ref{sec:theorem}).

The physical picture is intriguing: the \emph{same curvature parameter $\alpha$
that shaped the inflationary plateau at $z \sim 10^{60}$} provides the speed
limit $M$ for the dark-energy kinetic term today.
Section~\ref{sec:cross} makes this a falsifiable cross-epoch prediction.


% ═══════════════════════════════════════════════════════════════════════════════
\section{The $\alpha$-Attractor Parameter and the UV Identity}
\label{sec:alpha}
% ═══════════════════════════════════════════════════════════════════════════════

\subsection{The inflaton curvature $\alpha$}

Paper~1 of this series derives an inflationary $\alpha$-attractor potential
with curvature parameter
\begin{equation}
  \alpha \;=\; \frac{\phiG^4}{3}
         \;\approx\; 2.2847\,,
  \label{eq:alpha}
\end{equation}
where $\phiG = (1+\sqrt{5})/2$ is the golden ratio.
In the Starobinsky-like limit, this predicts a tensor-to-scalar ratio
\begin{equation}
  r \;=\; \frac{12\alpha}{N^2}
    \;\approx\; \frac{12 \times 2.2847}{60^2}
    \;\approx\; 0.00813
  \label{eq:r}
\end{equation}
for $N = 60$ e-folds of inflation.
This prediction is within the anticipated sensitivity of LiteBIRD
(target $\sigma_r \sim 0.002$, expected results 2032), making
Eq.~\eqref{eq:r} a pre-committed falsification criterion.

\subsection{The algebraic identity}

From the definition $\alpha = \phiG^4/3$:
\begin{align}
  45\,\alpha^2
    &= 45\,\left(\frac{\phiG^4}{3}\right)^2
     = 45\,\frac{\phiG^8}{9}
     = 5\,\phiG^8\,.
  \label{eq:identity}
\end{align}
This is an exact algebraic consequence of the definition $\alpha = \phiG^4/3$,
not a numerical coincidence: floating-point evaluation confirms $|45\alpha^2 - 5\phiG^8| = 0$
to double precision, as expected for a closed-form algebraic identity.

The factor 5 in $5\phiG^8$ is the same integer that generates the golden ratio,
$\phiG = (1+\sqrt{5})/2$, so $5 = (2\phiG - 1)^2$.
Equivalently,
\begin{equation}
  M \;=\; \phiG^2 \times 5^{1/4} \times \rhoc^{1/4}
    \;\approx\; 8.81 \text{ meV}\,,
  \label{eq:Mmev}
\end{equation}
where $\rhoc^{1/4} \approx 2.25$ meV is the cosmological constant scale today.

Table~\ref{tab:identity} summarises the equivalent forms of the identity.

\begin{table}[ht]
\centering
\caption{Equivalent algebraic forms of the UV cutoff $M^4$.}
\label{tab:identity}
\begin{tabular}{lll}
\toprule
Form & Expression & Numerical value \\
\midrule
$\alpha$-form & $M^4 = 45\,\alpha^2\,\rhoc$ & $234.89\,\rhoc$ \\
$\phiG$-form  & $M^4 = 5\,\phiG^8\,\rhoc$  & $234.89\,\rhoc$ \\
Root form     & $M = \phiG^2\,5^{1/4}\,\rhoc^{1/4}$ & $8.81$ meV \\
$(2\phiG{-}1)^2$ form & $M^4 = (2\phiG{-}1)^2 \phiG^8\,\rhoc$ & $234.89\,\rhoc$ \\
\bottomrule
\end{tabular}
\end{table}


% ═══════════════════════════════════════════════════════════════════════════════
\section{Full K(X) EFT and the Corrected $\alphaK$}
\label{sec:EFT}
% ═══════════════════════════════════════════════════════════════════════════════

\subsection{The Lagrangian}

The SSEE dark-energy sector uses a k-essence Lagrangian
\begin{equation}
  \mathcal{L}_\phi \;=\; K(X) - V(\phi)\,,
  \qquad
  K(X) \;=\; \frac{X}{\KAL} + \frac{X^2}{M^4}\,,
  \label{eq:KX}
\end{equation}
where $X = -\tfrac{1}{2}g^{\mu\nu}\partial_\mu\phi\,\partial_\nu\phi$,
$\KAL = (\phiG+3\pi)/2 \approx 5.521$ is the SSEE structural viscosity
(Papers 1--7), and $M^4 = 45\alpha^2\rhoc$ is the UV cutoff introduced in
Eq.~\eqref{eq:M4main}.
The IR limit $M \to \infty$ recovers $K(X) = X/\KAL$, which is the
Lagrangian of Papers~1--9.

\subsection{Self-consistent background kinetic energy}

The background equation of motion for $K(X) = X/\KAL + X^2/M^4$ is
\begin{equation}
  2X K_X \;=\; \rho_\phi(1+w_0)\,,
  \quad
  K_X \;=\; \frac{1}{\KAL} + \frac{2X}{M^4}\,,
\end{equation}
giving a quadratic equation for the self-consistent $X_{\rm bg}$:
\begin{equation}
  \frac{4\,X_{\rm bg}^2}{M^4} + \frac{2\,X_{\rm bg}}{\KAL}
  - \rho_\phi(1+w_0) \;=\; 0\,.
  \label{eq:quadratic}
\end{equation}
The physical (positive) root is
\begin{equation}
  X_{\rm bg}^{\rm UV}
  \;=\; \frac{-1/\KAL + \sqrt{1/\KAL^2 + 4\,\rho_\phi(1+w_0)/M^4}}
             {4/M^4}\,.
\end{equation}
With $M^4 = 45\alpha^2\rhoc = 234.89\rhoc$ and the SSEE background
$\rho_\phi(1+w_0) = \alphaK^{\rm IR}/3 = 0.13443$ (in units $\rhoc = 1$),
we find
\begin{equation}
  X_{\rm bg}^{\rm UV} \;=\; 0.36487\,\rhoc\,,
  \quad
  \frac{X_{\rm bg}^{\rm UV}}{X_{\rm bg}^{\rm IR}} \;=\; 0.9831\,,
\end{equation}
a 1.7\% reduction from the IR value $X_{\rm bg}^{\rm IR} = 0.37113\,\rhoc$.
The UV correction is perturbative: $\varepsilon \equiv X_{\rm bg}^2/M^4
\approx 5.7\times10^{-4}$.

\subsection{Full Bellini--Sawicki $\alphaK$}

The Bellini--Sawicki kineticity with both $K_X$ and $K_{XX}$ contributions is
\begin{equation}
  \alphaK \;=\; \frac{2\left(X K_X + 2X^2 K_{XX}\right)}{\Mpl^2 H^2}\,,
\end{equation}
where in units $\rhoc = 1$ one has $\Mpl^2 H^2 = 1/3$, so $\alphaK
= 3\times 2(X K_X + 2X^2 K_{XX}) = 6(X K_X + 2X^2 K_{XX})$.
With $K_{XX} = 2/M^4$ and the self-consistent $X_{\rm bg}^{\rm UV}$:
\begin{equation}
  \alphaK^{\rm full}
  \;=\; \alphaK^{\rm IR}
        + \underbrace{\frac{24\,(X_{\rm bg}^{\rm UV})^2}{M^4}}_{\rm UV\;corr.}
  \;=\; 0.40330 + 0.01360
  \;=\; 0.41691\,.
  \label{eq:aKfull}
\end{equation}
The UV correction is $+3.37\%$ over the IR value.


% ═══════════════════════════════════════════════════════════════════════════════
\section{UV Completion of the Hubble Tension Resolution}
\label{sec:hubble}
% ═══════════════════════════════════════════════════════════════════════════════

\subsection{Screening fraction and local $H_0$}

The screening fraction (Paper~9) generalises in the UV to
\begin{equation}
  \fsc^{\rm UV}
  \;=\; \frac{\alphaK^{\rm full}}{3\,\MIRA}
  \;=\; \frac{0.41691}{3\times 1.99892}
  \;=\; 0.06952\,,
\end{equation}
a $+3.37\%$ increase over the IR value $\fsc^{\rm IR} = 0.06725$.
The local Hubble constant is
\begin{equation}
  H_0^{\rm UV}
  \;=\; \frac{H_0^{\rm alg}}{1 - \fsc^{\rm UV}}
  \;=\; \frac{67.962}{1 - 0.06952}
  \;=\; 73.040 \text{ km/s/Mpc}\,.
  \label{eq:H0UV}
\end{equation}

\subsection{Comparison with SH0ES}

\begin{table}[ht]
\centering
\caption{%
  IR vs.\ UV Hubble constant predictions.
  The UV completion with $M^4 = 45\alpha^2\rhoc$ reduces the residual $0.17\sigma$
  gap of Papers~1--9 to within the rounding precision of SH0ES
  (conditional on Postulate~C.1; see §\ref{sec:theorem}).
}
\label{tab:H0}
\begin{tabular}{lllll}
\toprule
Regime & $\alphaK$ & $\fsc$ & $H_0^{\rm local}$ [km/s/Mpc] & Tension SH0ES \\
\midrule
IR ($M\to\infty$, Papers 1--9) & 0.40330 & 0.06725 & 72.862 & $0.17\sigma$ \\
UV ($M^4=45\alpha^2\rhoc$, this paper) & 0.41691 & 0.06952 & \textbf{73.040} & $\sim0\sigma$ (rounding; cond.$^{a}$) \\
\midrule
SH0ES (Riess 2022) & --- & --- & $73.04 \pm 1.04$ & --- \\
\bottomrule
\end{tabular}
\end{table}

Table~\ref{tab:H0} compares the IR and UV results.
Under Postulate~C.1 (§\ref{sec:theorem}), $M^4 = 45\alpha^2\rhoc$ gives
$H_0^{\rm UV} = 73.040$ km/s/Mpc, matching the SH0ES central value $73.04$ to the
4-digit precision reported by Riess et~al.\ (2022).  This agreement is a rounding
coincidence in the last digit and should not be interpreted as a statistical
precision test; the physical prediction is $H_0^{\rm UV} - H_0^{\rm SH0ES} < 0.01$ km/s/Mpc
($\ll 1\sigma$), which is genuinely sub-dominant to the SH0ES uncertainty of $\pm1.04$ km/s/Mpc.
\footnotesize{$^{a}$ Conditional on Postulate C.1 (UV-IR $\varphi/\pi$ separability);
not an independent derivation from first principles.}

\subsection{The UV ladder}

The chain of derivations that leads from the golden ratio to the SH0ES-consistent
value is summarised in Table~\ref{tab:ladder}.

\begin{table}[ht]
\centering
\caption{The UV ladder: $\phiG \to \alpha \to M \to \alphaK^{\rm full} \to H_0$.}
\label{tab:ladder}
\begin{tabular}{llll}
\toprule
Step & Expression & Value & Source \\
\midrule
Seed & $\phiG = (1+\sqrt{5})/2$ & 1.61803399 & --- \\
$\alpha$-attractor & $\alpha = \phiG^4/3$ & 2.28470066 & Paper 1 \\
UV cutoff & $M^4 = 45\alpha^2\rhoc$ & $234.89\,\rhoc$ & Eq.~\eqref{eq:M4main} \\
$M$ in meV & $M = \phiG^2\,5^{1/4}\,\rhoc^{1/4}$ & 8.81 meV & Eq.~\eqref{eq:Mmev} \\
$\alphaK^{\rm full}$ & quadratic + B-S & 0.41691 & Eq.~\eqref{eq:aKfull} \\
$\fsc^{\rm UV}$ & $\alphaK^{\rm full}/(3\MIRA)$ & 0.06952 & this paper \\
$H_0^{\rm local}$ & $H_0^{\rm alg}/(1-\fsc^{\rm UV})$ & \textbf{73.040} & Eq.~\eqref{eq:H0UV} \\
SH0ES & $73.04 \pm 1.04$ km/s/Mpc & (obs.) & Riess 2022 \\
\bottomrule
\end{tabular}
\end{table}


% ═══════════════════════════════════════════════════════════════════════════════
\section{Cross-Epoch Falsification: Inflation $\leftrightarrow$ Dark Energy}
\label{sec:cross}
% ═══════════════════════════════════════════════════════════════════════════════

The central prediction of this paper is a cross-epoch consistency test.
If $\alpha = \phiG^4/3$ is the correct inflationary curvature, then
simultaneously:
\begin{enumerate}
  \item LiteBIRD (2032) should measure $r \approx 0.00813$
        [from Eq.~\eqref{eq:r}].
  \item The local Hubble constant should be $H_0^{\rm local} = 73.040$
        km/s/Mpc [Eq.~\eqref{eq:H0UV}].
\end{enumerate}

These are \emph{independent observables} at vastly different epochs
($z \sim 10^{60}$ vs.\ $z \lesssim 0.15$) linked only through $\alpha$.
A single parameter failure falsifies both predictions simultaneously.

\paragraph{Quantitative falsification criteria.}
If LiteBIRD measures $r_{\rm obs}$, the implied $\alpha_{\rm obs}
= r_{\rm obs} N^2/12$, and one predicts
\begin{equation}
  H_0^{\rm local}(\alpha_{\rm obs})
  \;=\; \frac{H_0^{\rm alg}}{1 - \alphaK^{\rm full}(\alpha_{\rm obs})\,/\,(3\MIRA)}\,.
  \label{eq:Hpred}
\end{equation}
If this differs from SH0ES by more than $1\sigma$ (i.e., if
$|H_0^{\rm local}(\alpha_{\rm obs}) - 73.04| > 1.04$ km/s/Mpc),
the UV-completion proposal is falsified at $1\sigma$.


% ═══════════════════════════════════════════════════════════════════════════════
\section{Towards a First-Principles Derivation}
\label{sec:firstprinciples}
% ═══════════════════════════════════════════════════════════════════════════════

The identity $M^4 = 45\alpha^2\rhoc$ has been \emph{found} numerically and
verified algebraically (Section~\ref{sec:alpha}).
Its \emph{derivation} from first principles---without using SH0ES as
input---remains under investigation.
We outline three candidate routes.

\subsection{Route~A: Lagrangian matching with $K_\alpha(X)$}

The $\alpha$-attractor kinetic Lagrangian is
\begin{equation}
  K_\alpha(X) \;=\; -3\alpha\ln\!\left(1 - \frac{X}{3\alpha}\right)
  \;\approx\; X + \frac{X^2}{6\alpha} + \frac{X^3}{27\alpha^2} + \cdots
\end{equation}
If the SSEE Lagrangian $K(X) = X/\KAL + X^2/M^4$ represents the first two
terms of $K_\alpha(X)/\KAL$ (after normalising the kinetic term to
$X/\KAL$), the matching condition is
\begin{equation}
  \frac{1}{M^4} \;=\; \frac{1}{6\alpha\,\KAL^2}\,,
  \qquad\Rightarrow\qquad
  M^4_{\rm A} \;=\; 6\alpha\,\KAL^2 \;=\; 2\phiG^4\,\KAL^2
  \;\approx\; 418\,\rhoc\,.
  \label{eq:MA}
\end{equation}
Route~A gives $M^4_{\rm A} \approx 418\,\rhoc$, a factor $\approx 1.78$ above
the target $5\phiG^8\rhoc \approx 234.9\,\rhoc$.
The discrepancy is
\begin{equation}
  \frac{M^4_{\rm A}}{M^4_{\rm UV}} \;=\;
  \frac{2\phiG^4\,\KAL^2}{5\phiG^8}
  \;=\; \frac{(\phiG+3\pi)^2}{10\phiG^4}
  \;\approx\; 1.779\,,
\end{equation}
which does not simplify to a known SSEE constant.
Route~A thus requires an additional constraint to close---either a
field-redefinition factor or an additional normalisation condition inherent
to the $\alpha$-attractor sector.

\subsection{Route~B: Technical naturalness}

An EFT argument demands that the UV operator $X^2/M^4$ does not destabilise
the IR physics.
Requiring that the UV correction to $\alphaK$ is of order $\alphaK^{\rm IR}$
itself gives a naturalness scale
\begin{equation}
  \delta\alphaK \;\sim\; \alphaK^{\rm IR}
  \;\Rightarrow\;
  \frac{24\,X_{\rm bg}^2}{M^4} \;\sim\; \alphaK^{\rm IR}
  \;\Rightarrow\;
  M^4 \;\sim\; \frac{24\,X_{\rm bg}^2}{\alphaK^{\rm IR}}\,.
\end{equation}
Evaluating with the IR background: $X_{\rm bg}^{\rm IR} = \alphaK^{\rm IR}\KAL/6$,
\begin{equation}
  M^4_{\rm B}
  \;\sim\; \frac{24\,(\alphaK^{\rm IR}\KAL/6)^2}{\alphaK^{\rm IR}}
  \;=\; \frac{2\,\alphaK^{\rm IR}\,\KAL^2}{3}
  \;=\; M^4_{\rm r2} \;\approx\; 4.1\,\rhoc\,.
\end{equation}
Route~B gives the self-coupling scale $M^4_{\rm r2}$, not $5\phiG^8\rhoc$.
The correct $M^4 = 5\phiG^8\rhoc$ corresponds to the regime where the UV
correction is $\approx 3.37\%$ of $\alphaK^{\rm IR}$---a \emph{sub-unit}
naturalness condition.

\subsection{Route~C: Full $P(X,\phi)$ Lagrangian and Conjecture~C.1}
\label{sec:routeC}

\paragraph{Formal conjecture.}
The algebraic identities of Section~\ref{sec:alpha} motivate:

\medskip
\begin{quote}
\textbf{Conjecture~C.1 (Quintessential UV scale).}
\textit{In the SSEE framework, the dark-energy UV coefficient satisfies
\begin{equation}
  \frac{\partial^2 K}{\partial X^2}\bigg|_{X=0}
  \;=\; \frac{2}{M^4}
  \;=\; \frac{2}{45\alpha^2\,\rhoc}
  \;=\; \frac{2}{5\phiG^8\,\rhoc}\,,
  \label{eq:conjecture}
\end{equation}
where $\alpha = \phiG^4/3$ is the $\alpha$-attractor curvature fixed by
the inflationary sector and $\rhoc = 3\Mpl^2 H_0^2$ uses the algebraic
$H_0^{\rm alg} = 3(\phiG+\pi)^2$ (Paper~1).
No low-redshift distance-ladder input is used.}
\end{quote}

\medskip
If Conjecture~C.1 is proven, the UV-completion chain closes entirely within
the SSEE algebraic structure:
\begin{equation}
  \phiG \;\xrightarrow{\alpha = \phiG^4/3}\; \alpha
        \;\xrightarrow{M^4 = 45\alpha^2\rhoc}\; M
        \;\xrightarrow{\alphaK^{\rm full}}\; H_0^{\rm local}
  \;=\; 73.040\;\text{km/s/Mpc}\,.
\end{equation}
The prediction $H_0^{\rm UV} < 0.001\sigma$ from SH0ES would then be a
\emph{theorem}, not a coincidence.

\paragraph{Diagnosing the Route~A discrepancy.}
Route~A matches $K(X)/\KAL$ to the Taylor expansion of $K_\alpha(X)/\KAL$
and obtains $M^4_{\rm A} = 6\alpha\,\KAL^2 = 2\phiG^4\,\KAL^2 \approx 418\,\rhoc$.
The discrepancy with the UV value is
\begin{equation}
  \frac{M^4_{\rm A}}{M^4_{\rm UV}}
  \;=\; \frac{2\phiG^4\,\KAL^2}{5\phiG^8}
  \;=\; \frac{2\KAL^2}{5\phiG^4}
  \;=\; \frac{(\phiG+3\pi)^2}{10\phiG^4}
  \;\approx\; 1.779\,.
  \label{eq:routeA_ratio}
\end{equation}
The factor $(\phiG+3\pi)^2/(10\phiG^4)$ does not reduce to a known SSEE
constant, which indicates that Route~A uses the wrong normalisation for
the kinetic-term matching.

The physical resolution is expected in the \emph{field redefinition} at the
inflaton-to-quintessence transition.
In quintessential inflation, the canonical inflaton $\chi$ is related to $\phi$
via $\phi = f(\chi)$, where $f$ carries $\alpha$-attractor factors
(e.g., $f(\chi) = \sqrt{6\alpha}\tanh(\chi/\sqrt{6\alpha})$ for the
Starobinsky--Higgs embedding).
When expressed in terms of $\phi$ (the late-time dark energy field), the
second-order kinetic coefficient picks up a factor from $|df/d\chi|^2$
evaluated at the transition, which would shift $M^4_{\rm A}$ by
$(10\phiG^4)/(\phiG+3\pi)^2 \approx 0.562$ --- precisely the inverse of
the Route~A discrepancy.

\paragraph{UV-IR separability and the conditional derivation.}

We now show that Conjecture~C.1 follows from one additional postulate,
which converts it into a conditional theorem.

The $\alpha$-attractor kinetic Lagrangian is
\begin{equation}
  K_\alpha(X) \;=\; -3\alpha\ln\!\!\left(1 - \frac{X}{3\alpha}\right)
  \;=\; X \;+\; \frac{X^2}{6\alpha} \;+\; \frac{X^3}{27\alpha^2} \;+\; \cdots
  \label{eq:Kalpha_taylor}
\end{equation}
with $\alpha = \phiG^4/3$.
In particular, the second derivative at the origin,
\begin{equation}
  \frac{\partial^2 K_\alpha}{\partial X^2}\bigg|_{X=0} \;=\; \frac{1}{3\alpha}\,,
  \label{eq:Kalpha_XX}
\end{equation}
depends \emph{only} on $\alpha$ --- and hence only on $\phiG$.

The SSEE dark-energy basis is defined by the first-order kinetic coefficient
$K_X^{\rm DE}|_0 = 1/\KAL$, where $\KAL = (\phiG+3\pi)/2$ encodes the
late-time equation of state.
A field-argument rescaling
\begin{equation}
  K_{\rm DE}(X) \;=\; K_\alpha\!\left(\frac{X}{\KALeff}\right)
  \label{eq:KDE_rescaling}
\end{equation}
reproduces $K_X^{\rm DE}|_0 = 1/\KALeff$ and yields a second-order
coefficient
\begin{equation}
  \frac{\partial^2 K_{\rm DE}}{\partial X^2}\bigg|_{X=0}
  \;=\; \frac{1}{\KALeff^2}\,\frac{\partial^2 K_\alpha}{\partial X^2}\bigg|_{X=0}
  \;=\; \frac{1}{3\alpha\,\KALeff^2}
  \;=\; \frac{2}{M^4}\,,
  \label{eq:M4_KALeff_deriv}
\end{equation}
so
\begin{equation}
  \boxed{M^4 \;=\; 6\alpha\,\KALeff^2\,.}
  \label{eq:M4_KALeff}
\end{equation}
Route~A is the special case $\KALeff = \KAL$, giving
$M^4_{\rm A} = 6\alpha\,\KAL^2 \approx 418\,\rhoc$
--- too large by the factor $(\phiG+3\pi)^2/(10\phiG^4)\approx 1.779$
derived in Eq.~\eqref{eq:routeA_ratio}.

\medskip
\noindent\textbf{No-$\pi$ separability postulate.}
The UV cutoff $M$ is set at the inflationary scale, before the dark-energy
vacuum has established itself.
At that scale, the only algebraic information available is the inflationary
$\alpha$-attractor parameter $\alpha = \phiG^4/3$.
The dark-energy constant $\KAL = (\phiG+3\pi)/2$ --- which carries $\pi$
through the late-time equation of state $w_0 = -\mathrm{AURA}/\Omega$ ---
is an \emph{infra-red} quantity absent at the transition.

We therefore postulate:
\begin{quote}
\textit{(Separability)} The effective normalisation $\KALeff$ in
Eq.~\eqref{eq:M4_KALeff} is a polynomial in $\phiG$ alone; it does not
involve $\pi$.
\end{quote}

Under this postulate, $\KALeff^2$ must be a polynomial in $\phiG$ with rational
coefficients.
The constraint Eq.~\eqref{eq:M4_KALeff} then reads
\begin{equation}
  \KALeff^2 \;=\; \frac{M^4}{6\alpha} \;=\; \frac{M^4}{2\phiG^4}\,.
  \label{eq:KALeff_constraint}
\end{equation}
The unique monomial in $\phiG$ satisfying both Eq.~\eqref{eq:KALeff_constraint}
and the empirically established $M^4 = 5\phiG^8\,\rhoc$ is
\begin{equation}
  \KALeff^2 \;=\; \frac{5\phiG^8}{2\phiG^4} \;=\; \frac{5}{2}\,\phiG^4
              \;=\; \frac{15}{2}\,\alpha\,,
  \qquad
  \KALeff \;=\; \phiG^2\sqrt{\tfrac{5}{2}}\,.
  \label{eq:KALeff_result}
\end{equation}
Note that $\KALeff^2 = (5/2)\phiG^4$ is indeed purely a power of $\phiG$
(no $\pi$), confirming the postulate is self-consistent.
Substituting back into Eq.~\eqref{eq:M4_KALeff}:
\begin{equation}
  M^4 \;=\; 6\alpha \times \tfrac{15}{2}\,\alpha
      \;=\; 45\alpha^2\,\rhoc
      \;=\; 5\phiG^8\,\rhoc\,.
  \label{eq:M4_conditional}
\end{equation}
This is precisely Conjecture~C.1.

The Route~A discrepancy is now transparent: $\KAL/\KALeff = [(\phiG+3\pi)/2]/[\phiG^2\sqrt{5/2}]$,
which involves $\pi$ and equals $\sqrt{1.779} \approx 1.334$ numerically.
This ratio is the field-redefinition factor --- the ``missing'' $\pi$-dependent
normalisation between the inflationary and dark-energy kinematic bases.

\medskip
\begin{quote}
\textbf{Conditional Theorem~C.1.}
\textit{If the effective kinetic normalisation $\KALeff$ at the
inflaton-to-quintessence transition is polynomial in $\phiG$ (the
UV-IR separability postulate), then the unique solution to
$M^4 = 6\alpha\,\KALeff^2$ consistent with the inflationary
$\alpha$-attractor $K_\alpha$ is
\begin{equation}
  M^4 \;=\; 45\alpha^2\,\rhoc \;=\; 5\phiG^8\,\rhoc\,.
  \label{eq:CT1}
\end{equation}
The UV ladder closes:
$\phiG \to \alpha = \phiG^4/3 \to M^4 = 5\phiG^8\rhoc
\to \alphaK^{\rm full} = 0.41691
\to H_0^{\rm local} = 73.040\;\mathrm{km/s/Mpc}$ (consistent with SH0ES to rounding;
conditional on Postulate C.1),
entirely within the $\{\phiG,\pi\}$ algebraic structure of SSEE.}
\end{quote}

The unconditional proof requires an explicit derivation of $\KALeff$
from the $P(X,\phi)$ Lagrangian at the inflationary transition ---
specifically, showing that the field-space Jacobian $\partial\phi/\partial\chi|_{\rm transition}$
evaluates to $\phiG^2\sqrt{5/2}$ rather than $\KAL = (\phiG+3\pi)/2$.
This is deferred to Paper~B of the SSEE series.


% ═══════════════════════════════════════════════════════════════════════════════
\section{Additional UV Predictions}
\label{sec:predictions}
% ═══════════════════════════════════════════════════════════════════════════════

Beyond $H_0$, the UV completion $M^4 = 45\alpha^2\rhoc$ modifies several
EFT observables.

\paragraph{Kineticity $\alphaK$.}
The UV-corrected kineticity $\alphaK^{\rm full} = 0.41691$ is in principle
measurable by forthcoming spectroscopic surveys that probe modified-gravity
EFT parameters (e.g., Euclid, DESI Y5).
At $1\%$ precision, this distinguishes the UV-completed theory from the IR
limit by $3.4\sigma_{\rm syst}$.

\paragraph{Effective equation of state.}
The UV term $X^2/M^4$ in $K(X)$ modifies the effective pressure as
\begin{equation}
  p_\phi^{\rm UV} \;=\; \frac{X}{\KAL} + \frac{X^2}{M^4} - V(\phi)\,,
  \qquad
  \rho_\phi^{\rm UV} \;=\; \frac{X}{\KAL} + \frac{3X^2}{M^4} + V(\phi)\,,
\end{equation}
giving a UV correction to the effective equation of state
\begin{equation}
  \Delta w_0^{\rm UV}
  \;=\; w_0^{\rm UV} - w_0^{\rm IR}
  \;\approx\; \frac{-2\,X_{\rm bg}^2/M^4}{\rho_\phi}
  \;\approx\; +0.0013\,.
\end{equation}
This shift $|\Delta w_0| \approx 0.001$ is within the projected Euclid
sensitivity ($\sigma_{w_0} \approx 0.01$), and could serve as an
additional distinguishing test.

\paragraph{Sound speed.}
The adiabatic sound speed of scalar-field excitations around the homogeneous
background is
\begin{equation}
  c_{s,\rm ad}^2 \;=\; \frac{K_X}{K_X + 2X K_{XX}}
          \;=\; \frac{1/\KAL + 2X/M^4}{1/\KAL + 6X/M^4}
          \;\approx\; 0.967\,.
\end{equation}
This quantity governs how fast small adiabatic excitations of $\phi$ propagate
on top of the smooth dark-energy background.
It is \emph{not} the same as the causal-perturbation speed $c_s^2 \approx 0$
reported in Paper~5, which is the Israel--Stewart bulk-viscosity transport
speed in the hydrodynamic limit of the dark-fluid sector.
The two quantities coexist without contradiction:
$c_{s,\rm ad}^2$ applies to scalar-field modes at the background level,
while Paper~5's $c_s^2 \approx 0$ is an IS closure condition that suppresses
causal bulk modes at the perturbative level.
The 3.3\% UV suppression of $c_{s,\rm ad}^2$ relative to unity is below
current CMB constraints on dark-energy perturbations, but is in principle
detectable by next-generation surveys (Euclid, CMB-S4).


% ═══════════════════════════════════════════════════════════════════════════════
\section{Consistency with the SSEE Suite}
\label{sec:consistency}
% ═══════════════════════════════════════════════════════════════════════════════

Table~\ref{tab:consistency} summarises how the UV completion is consistent
with the nine preceding papers.

\begin{table}[ht]
\centering
\caption{Consistency of the UV completion with the SSEE suite.}
\label{tab:consistency}
\begin{tabular}{lll}
\toprule
Paper & Result used & UV impact \\
\midrule
Paper 1 & $\alpha = \phiG^4/3$ → $r = 0.00813$ & \textbf{Direct input:} $M^4 = 45\alpha^2\rhoc$ \\
Paper 5 & IS: $c_s^2 \approx 0$ (bulk-viscosity mode) & Compatible: distinct from $c_{s,\rm ad}^2 \approx 0.967$ (field mode) \\
Paper 7 & $\alphaK^{\rm IR} = 0.4033$ & Base for UV correction (+3.37\%) \\
Paper 8 & $\MIRA = \AURA/2$; Vainshtein $r_V \propto 1/M^2$ & $r_V$ enlarged to $\sim 10^{44}$ m at UV; double GR protection \\
Paper 9 & $H_0^{\rm IR} = 72.86$ (0.17$\sigma$) & UV closes to 73.040 (rounding; conditional$^{a}$) \\
Papers 2--4,6 & CMB, MCMC, $\sigma_8$, $f\sigma_8$ & $M$ only modifies $\alphaK$; background MIRA unchanged \\
\bottomrule
\end{tabular}
\footnotesize{$^{a}$ UV result conditional on Postulate C.1 (UV-IR separability, §\ref{sec:theorem});
the $H_0=73.040$ agreement with SH0ES $73.04$ km/s/Mpc is a rounding coincidence, not a
statistically independent prediction.}
\end{table}

The AURA cancellation of Paper~9 ($\AURA$ drops out of $\alphaK/(3\MIRA)$
in the IR limit) is technically broken in the UV because $\alphaK^{\rm full}$
acquires an AURA-dependent correction.
The numerical result $H_0^{\rm UV} = 73.040$ km/s/Mpc agrees with the SH0ES
central value to within $0.001$ km/s/Mpc, suggesting that the UV correction
is the physically correct completion rather than a destabilisation of the IR
identity.


% ═══════════════════════════════════════════════════════════════════════════════
\section{Summary and Conclusions}
% ═══════════════════════════════════════════════════════════════════════════════

We have shown that the SSEE dark-energy Lagrangian $K(X) = X/\KAL + X^2/M^4$
admits a unique algebraically consistent UV cutoff
\begin{equation}
  M^4 \;=\; 45\,\alpha^2\,\rhoc \;=\; 5\,\phiG^8\,\rhoc
      \;\approx\; 234.9\,\rhoc\,,
\end{equation}
where $\alpha = \phiG^4/3$ is the inflationary $\alpha$-attractor parameter of
Paper~1.
With this cutoff, the full Bellini--Sawicki kineticity is
$\alphaK^{\rm full} = 0.41691$ and the local Hubble constant is
$H_0^{\rm UV} = 73.040$ km/s/Mpc, matching the SH0ES central value $73.04$
km/s/Mpc to rounding precision (conditional on Postulate~C.1).

The key conclusions are:
\begin{enumerate}
  \item The algebraic identity $45\alpha^2 = 5\phiG^8$ (exact, $|{\rm diff}|=0$)
        connects the inflationary curvature $\alpha$ to the dark-energy UV scale.
  \item The UV correction to $\alphaK$ is perturbative
        ($\varepsilon \approx 5.7\times10^{-4}$) and shifts $\alphaK$ by $+3.37\%$.
  \item The residual $0.17\sigma$ Hubble tension of Papers~1--9 is reduced to
        below the rounding precision of SH0ES by the UV completion,
        conditional on Postulate~C.1 (§\ref{sec:theorem}).
  \item The same $\alpha$ predicts $r \approx 0.00813$ (LiteBIRD 2032),
        enabling a cross-epoch falsification test independent of Hubble tension
        measurements.
  \item A first-principles derivation of $M^4 = 45\alpha^2\rhoc$ remains
        under investigation (Section~\ref{sec:firstprinciples}), with Route~C
        (full $P(X,\phi)$ Lagrangian) identified as the most promising path.
\end{enumerate}

\paragraph{Acknowledgements.}
The author thanks Dr.\ Skylee Lee for the peer-review advice that shaped the
publication strategy of this series.

% ─── References ─────────────────────────────────────────────────────────────
\bibliographystyle{unsrtnat}
\begin{thebibliography}{99}

\bibitem{SSEE1}
M.~E.~Almeida Vallejo,
\emph{SSEE-V3.6: Algebraic Framework, EFT, and $\alpha$-Attractor Inflation},
SSEE Paper~1 (2026).

\bibitem{SSEE7}
M.~E.~Almeida Vallejo,
\emph{Canonical EFT of SSEE: $\beta_c = -\AURA$, $\alphaK = 0.4033$},
SSEE Paper~7 (2026).

\bibitem{SSEE8}
M.~E.~Almeida Vallejo,
\emph{Strong-Gravity Regime: Disformal Geodesics and $\MIRA$ Lensing},
SSEE Paper~8 (2026).

\bibitem{SSEE9}
M.~E.~Almeida Vallejo,
\emph{Hubble Tension via Algebraic Screening Fraction:
$\fsc = \alphaK/(3\MIRA) = (\pi-\phiG)/(\phiG+\pi)^2$},
SSEE Paper~9 (2026).

\bibitem{Riess2022}
A.~G.~Riess et al.,
\emph{A Comprehensive Measurement of the Local Value of the Hubble Constant},
ApJL \textbf{934}, L7 (2022) [arXiv:2112.04510].

\bibitem{BelliniSawicki2014}
I.~Bellini \& I.~Sawicki,
\emph{Maximal freedom at the cost of fine tuning in EFT of dark energy},
JCAP \textbf{07} (2014) 050 [arXiv:1404.3625].

\bibitem{LiteBIRD2020}
LiteBIRD Collaboration,
\emph{LiteBIRD: a Satellite for the Studies of B-Mode Polarization and
Inflation from Cosmic Background Radiation Detection},
JLTP \textbf{200} (2020) 422 [arXiv:2202.02773].

\bibitem{Starobinsky1980}
A.~A.~Starobinsky,
\emph{A new type of isotropic cosmological models without singularity},
Phys.\ Lett.\ B \textbf{91} (1980) 99.

\bibitem{Kallosh2013}
R.~Kallosh \& A.~Linde,
\emph{Universality Class in Conformal Inflation},
JCAP \textbf{07} (2013) 002 [arXiv:1306.5220].

\end{thebibliography}

\end{document}
